\documentclass[a4paper,12pt,twoside,brazil]{abntex2}

\usepackage[T1]{fontenc}
\usepackage[utf8]{inputenc}
\usepackage[brazil]{babel}
\usepackage{geometry}
\usepackage{lmodern}
\usepackage{amsmath}
\usepackage{setspace}
\usepackage{hyperref}

\geometry{
    top=3cm,
    bottom=2cm,
    left=3cm,
    right=2cm
}

\setlength{\parindent}{1.25cm}
\setlength{\parskip}{0pt}
\OnehalfSpacing

\hypersetup{
    colorlinks=true,
    linkcolor=black,
    citecolor=black,
    urlcolor=black,
    pdftitle={Endogeneização da Natureza na Teoria da Produção},
    pdfauthor={Fernando Souza da Vieira}
}

\begin{document}

\begin{center}
    {\Large \bfseries Capital natural, instituições e alocação de políticas ambientais -- Tecnologia e escolhas privada e pública em Mato Grosso }\par\vspace{0.8cm}
    Fernando Souza da Vieira
\end{center}

\vspace{0.8cm}

\noindent\textbf{Resumo}

Lorem ipsum dolor sit amet, consectetur adipiscing elit. Sed non risus. Suspendisse lectus tortor, dignissim sit amet, adipiscing nec, ultricies sed, dolor.

Cras elementum ultrices diam. Maecenas ligula massa, varius a, semper congue, euismod non, mi.



\section{Introdução}
\label{sec:introducao}

% ECONOMIA --> MEIO AMBIENTE --> ECONOMIA
% Retroalimentação: decisões econômicas afetam o estoque ambiental, que pode alterar as condições produtivas e os retornos econômicos.

% HIPÓTESES
% ---
% HG: ECONOMIA E MEIO AMBIENTE FORMAM UM SISTEMA INTERDEPENDENTE
% Retornos + instituições --> conversão/conservação --> estoque de vegetação
% Estoque de vegetação × choques climáticos --> produtividade --> retornos econômicos

% H1: RETORNOS E INSTITUIÇÕES AFETAM CONVERSÃO/CONSERVAÇÃO
% Retorno econômico + instituições --> decisão de uso da terra --> estoque de vegetação
% Heterogeneidade territorial/biomas pode alterar incentivos e respostas dos agentes.

% H2: VEGETAÇÃO MODIFICA A RESPOSTA DA PRODUTIVIDADE A CHOQUES CLIMÁTICOS
% Choque climático (exógeno) × estoque de vegetação --> produtividade
% Heterogeneidade entre biomas pode alterar a magnitude dessa relação.

% EXTENSÃO POSSÍVEL:
% Informação conjunta de H1 e H2 pode subsidiar desenho e focalização territorial de políticas de conservação. Não constitui, por ora, hipótese independente.

% EIXOS CONCEITUAIS QUE CONECTAM O ARTIGO À TESE:
% ---
% ENDOGENEIDADE AMBIENTAL
% O estado ambiental não é apenas resultado da atividade econômica:
% ele também pode alterar resultados produtivos e incentivos econômicos.

% GOVERNANÇA TERRITORIAL
% Instituições, enforcement e incentivos operam sobre agentes e territórios heterogêneos, produzindo respostas distintas de conversão e conservação.

% INTERTEMPORALIDADE
% Decisões presentes de uso da terra alteram o estoque ambiental relevante para períodos posteriores e podem, por essa via, modificar condições futuras de produção. Por ora, é mecanismo conceitual, não hipótese testada diretamente.

% ------------------------------------
%%% (I) CONTEXTO
% ------------------------------------

%% Conflito: preservar ou converter.
A expansão agropecuária costuma ser contextualizada como uma escolha entre conservar a vegetação nativa para fins de sustentabilidade ou converter a terra para obter maior retorno econômico {CITAÇÃO}. Essa representação, contudo, simplifica uma decisão que não depende apenas da aptidão agronômica da terra ou da rentabilidade esperada da produção. A conversão da vegetação é uma decisão privada realizada sob preços, custos de produção, restrições de crédito, expectativas, liquidez, direitos de propriedade, regras ambientais, \textit{enforcement} e probabilidade esperada de detecção e sanção {CITAÇÃO}. 

Assim, territórios com retornos agrícolas esperados semelhantes podem apresentar decisões distintas de uso da terra quando diferem na configuração ou na efetividade das instituições. De modo similar, uma mesma regra formal pode produzir respostas diferentes conforme a atividade econômica, o custo de oportunidade da terra e as condições locais de fiscalização {CITAÇÃO}. Sob essa perspectiva, a conversão ou conservação da vegetação pode ser tratada como resultado endógeno da interação entre os retornos econômicos dos usos alternativos da terra e os incentivos, custos e riscos associados ao ambiente institucional {CITAÇÃO}.

% Ponte para H1
Mudanças nos preços relativos, na rentabilidade agrícola ou nas condições de acesso aos mercados podem aumentar o retorno esperado da conversão, enquanto fiscalização, sanções, restrições legais ou outros mecanismos institucionais podem alterar a resposta dos agentes a esses incentivos {CITAÇÃO}. A intensidade dessa resposta também pode variar entre territórios, atividades e biomas {CITAÇÃO}. Portanto, compreender a evolução do estoque de vegetação requer analisar não apenas os retornos econômicos associados ao uso da terra, mas também as instituições sob as quais esses retornos são percebidos e transformados em decisões de conversão ou conservação.

% Ponte para H2
Entretanto, o estoque ambiental afetado por essas decisões não constitui apenas uma consequência ambiental externa ao processo produtivo. A vegetação nativa interfere nas condições biofísicas que sustentam a produção e pode, por essa via, afetar os próprios resultados econômicos que influenciam o uso da terra {CITAÇÃO}. Ao influenciar ou estar associada a processos biofísicos relacionados à umidade, precipitação, temperatura e disponibilidade hídrica, por exemplo, o estoque de vegetação pode modificar as condições biofísicas locais e, por essa via, a resposta da agricultura aos choques climáticos {CITAÇÃO}. Nesse sentido, a vegetação não constitui apenas um ativo ambiental externo à firma, mas pode também atuar como condicionante da tecnologia produtiva, ainda que não seja um fator convencional diretamente controlado pelo produtor ou perfeitamente substituível por capital e outros insumos {CITAÇÃO}.

Essa distinção é particularmente importante porque variações climáticas de curto prazo podem fornecer uma fonte de variação plausivelmente exógena para investigar se a resposta da produtividade aos choques depende do estoque de vegetação previamente existente. Se áreas com diferentes estoques prévios de vegetação respondem de maneira distinta a choques comparáveis de precipitação, temperatura ou déficit hídrico, então o estoque ambiental pode modificar o efeito desses choques sobre a produtividade agrícola {CITAÇÃO}. A questão deixa, portanto, de ser apenas se a vegetação está correlacionada com maior ou menor produtividade e passa a ser se ela altera a sensibilidade da produção a perturbações externas, e se essa relação apresenta magnitude economicamente relevante e heterogeneidade entre biomas.

% Justificativa para a integração das duas hipóteses
Os dois mecanismos formam, em conjunto, uma relação de codependência entre economia e meio ambiente. De um lado, retornos econômicos e instituições condicionam a decisão de converter ou conservar a vegetação e, consequentemente, a evolução do estoque ambiental. De outro, esse estoque pode modificar as condições biofísicas da produção e a resposta da atividade agrícola a choques climáticos, afetando a produtividade e os retornos econômicos associados ao uso da terra {CITAÇÃO}. Assim, decisões econômicas alteram o estado ambiental, enquanto o estado ambiental pode alterar os resultados econômicos que fundamentam decisões subsequentes.

Essa retroalimentação pode ser especialmente relevante quando seus componentes são espacialmente heterogêneos. Os retornos da conversão, os custos institucionais enfrentados pelos agentes e a contribuição produtiva da vegetação podem variar entre atividades, territórios e biomas {CITAÇÃO}. Consequentemente, políticas que tratam a conservação apenas como restrição à atividade econômica, ou que consideram separadamente os incentivos à conversão e os efeitos produtivos do estoque ambiental, podem deixar de capturar parte relevante da relação entre produção e conservação. Surge, então, uma questão econômica central: em que medida a interação entre incentivos privados, instituições e contribuição produtiva do estoque ambiental pode informar o desenho e a focalização territorial das políticas de conservação?

%% Mato Grosso neste contexto.
Mato Grosso constitui um contexto especialmente informativo para estudar essas relações. O estado combina produção agrícola e pecuária em larga escala, áreas submetidas a diferentes pressões de conversão e territórios pertencentes à Amazônia, ao Cerrado e ao Pantanal. Essa configuração reúne, dentro de um mesmo ambiente estadual, diferenças relevantes de vegetação, clima, disponibilidade hídrica, sistemas produtivos e custo de oportunidade da terra. Os três biomas, portanto, não representam apenas distintas formações ecológicas. Eles também constituem contextos econômicos nos quais podem variar tanto a contribuição do estoque de vegetação para as condições produtivas quanto os incentivos privados associados à conversão ou à conservação da terra {CITAÇÃO}.

Essa heterogeneidade permite investigar empiricamente os dois mecanismos discutidos anteriormente. Do lado da decisão de uso da terra, a mesma heterogeneidade territorial permite examinar se retornos econômicos e instituições produzem respostas distintas de conversão entre atividades e biomas. Nesse caso, não basta saber se uma regra ambiental ou uma alteração na rentabilidade afeta, em média, o desmatamento ou a conversão. Importa também verificar se a intensidade dessas respostas depende das condições econômicas e territoriais sobre as quais os incentivos incidem. Assim, a heterogeneidade entre Amazônia, Cerrado e Pantanal interessa à pesquisa não apenas como controle espacial, mas como possível componente dos próprios mecanismos econômicos estudados. Do lado da produção, interessa saber se o estoque de vegetação modifica a resposta da produtividade aos choques climáticos e se essa relação difere entre biomas. O objeto de interesse não é, portanto, apenas verificar se a vegetação está associada ao nível médio de produtividade, mas também se ela altera a intensidade com que condições climáticas adversas afetam os resultados produtivos.



% ------------------------------------
%%% (II) PROBLEMA
% ------------------------------------

Diante desse quadro, o problema central deste artigo consiste em investigar como os incentivos econômicos e institucionais associados ao uso da terra se relacionam com a conversão ou manutenção da vegetação e se o estoque de vegetação previamente existente modifica a resposta da produtividade agrícola a choques climáticos em contextos ecológica e economicamente heterogêneos. A articulação desses dois mecanismos permite examinar uma relação de interdependência na qual decisões econômicas afetam o estado ambiental, enquanto o estado ambiental pode, por sua vez, condicionar resultados produtivos relevantes para decisões econômicas subsequentes.

O \textbf{problema} do artigo é responder ao seguinte questionamento: "Como os retornos econômicos e as instituições condicionam as decisões de conversão ou manutenção da vegetação e em que medida o estoque de vegetação nativa modifica a resposta da produtividade agrícola a choques climáticos em territórios ecologicamente e economicamente heterogêneos?"

Essa pergunta desdobra-se em duas dimensões empiricamente relacionadas: (i) como retornos econômicos, custos, riscos e incentivos institucionais se relacionam com as decisões de conversão ou manutenção da vegetação; e (ii) se o estoque de vegetação previamente existente modifica, em magnitude economicamente relevante, o efeito de choques climáticos sobre a produtividade agrícola e se essa relação varia entre biomas. Em conjunto, essas duas dimensões buscam fornecer informação sobre como incentivos privados, instituições e condições ambientais interagem na relação entre produção e conservação, podendo subsidiar o desenho e a focalização territorial de políticas ambientais.

% ------------------------------------
%%% (III) OBJETIVOS
% ------------------------------------

O \textbf{objetivo geral} deste artigo é investigar como retornos econômicos e instituições se relacionam com as decisões de conversão ou manutenção da vegetação e se o estoque de vegetação previamente existente modifica a resposta da produtividade agrícola a choques climáticos em contextos territorialmente heterogêneos.

De forma específica, o artigo busca: (i) examinar como retornos econômicos e incentivos institucionais se relacionam com as decisões de conversão ou manutenção da vegetação; (ii) verificar se o estoque de vegetação previamente existente modifica a resposta da produtividade agrícola a choques climáticos; e (iii) avaliar se esses mecanismos apresentam heterogeneidade entre biomas.

% ------------------------------------
%%% (IV) HIPÓTESES
% ------------------------------------

Dito isso, a análise parte da seguinte \textbf{hipótese geral (HG)}: retornos econômicos e instituições condicionam decisões de uso da terra que alteram o estoque ambiental, enquanto esse estoque pode, por sua vez, condicionar resultados produtivos relevantes para decisões econômicas subsequentes.

Essa hipótese geral é examinada por meio de duas hipóteses específicas.

\medskip

\noindent\textbf{(H1) Retornos, instituições e conversão.}

\emph{A resposta da conversão da vegetação aos retornos econômicos dos usos alternativos da terra é condicionada pelos incentivos, custos e riscos associados ao ambiente institucional, podendo variar entre atividades, territórios e biomas.}

\medskip

\noindent\textbf{(H2) Vegetação, clima e produtividade.}

\emph{O estoque de vegetação previamente existente modifica, em magnitude economicamente relevante, a resposta da produtividade agrícola aos choques climáticos, podendo essa relação variar entre biomas.}

% ------------------------------------
%%% (V) METODOLOGIA
% ------------------------------------

% detalhar dados/variáveis, fontes e período.
As hipóteses serão examinadas com dados em painel para Mato Grosso. H1 relacionará a conversão da vegetação aos retornos econômicos dos usos alternativos da terra, ao ambiente institucional e à interação entre essas dimensões. H2 relacionará a produtividade agrícola aos choques climáticos, ao estoque prévio de vegetação e à interação entre essas variáveis.

As especificações considerarão a ordenação temporal compatível com cada mecanismo, efeitos fixos territoriais e temporais e controles pertinentes a cada hipótese. Dependência espacial, heterogeneidade entre biomas e especificações alternativas serão examinadas por meio de testes e análises de robustez. A interpretação causal dos resultados dependerá das condições de identificação sustentadas pelos dados.

% ------------------------------------
%%% (VI) DIVISÃO DO TRABALHO
% ------------------------------------

Além desta introdução, o artigo apresenta o contexto econômico do problema e a literatura relacionada aos mecanismos investigados, desenvolve o referencial teórico, descreve a estratégia empírica e os dados, reporta os resultados, discute suas implicações e limitações e, por fim, apresenta as considerações finais.

\section{Contexto do Problema}
% (==>) É assim que esse mecanismo se manifesta concretamente no Brasil e em Mato Grosso.

% Aqui entra muito mais realidade institucional e empírica:
% expansão da soja/pecuária;
% estrutura fundiária;
% Código Florestal;
% fiscalização;
% Amazônia/Cerrado/Pantanal;
% desmatamento;
% precipitação;
% sistemas produtivos;
% instituições brasileiras.

% 2. Contexto econômico e institucional do problema
% 2.1. Produção agropecuária e conversão da vegetação
% Retornos, terra, soja, pecuária, expansão agrícola.

% 2.2. Instituições ambientais e decisões de uso da terra
% Código Florestal, fiscalização, sanções, regularização etc.

% 2.3. Vegetação, clima e produção agropecuária
% Mecanismos biofísicos.

% 2.4. Mato Grosso como ambiente empírico
% Amazônia, Cerrado, Pantanal e heterogeneidade territorial.

Lorem ipsum dolor sit amet, consectetur adipiscing elit. Sed do eiusmod tempor incididunt ut labore et dolore magna aliqua.

Ut enim ad minim veniam, quis nostrud exercitation ullamco laboris nisi ut aliquip ex ea commodo consequat.

\section{Revisão da Literatura}
% (==>) O que sabemos cientificamente sobre esse mecanismo, como sabemos, onde estão as controvérsias e o que ainda não sabemos?

% 3. Literatura relacionada

% 3.1. Retornos econômicos, instituições e conversão da vegetação — H1
% O que a literatura econômica já sabe.

% 3.2. Estoque de vegetação, clima e produtividade — H2
% O que as literaturas econômica, agronômica e climática já sabem.

% 3.3. A lacuna: integração dos dois mecanismos:
% economia e instituições --> E --> condições produtivas --> resultados econômmicos
% e sua heterogeneidade territorial.

Lorem ipsum dolor sit amet, consectetur adipiscing elit. Sed do eiusmod tempor incididunt ut labore et dolore magna aliqua.

Ut enim ad minim veniam, quis nostrud exercitation ullamco laboris nisi ut aliquip ex ea commodo consequat.

\section{Modelo Teórico}

Lorem ipsum dolor sit amet, consectetur adipiscing elit. Sed do eiusmod tempor incididunt ut labore et dolore magna aliqua.

Ut enim ad minim veniam, quis nostrud exercitation ullamco laboris nisi ut aliquip ex ea commodo consequat.

\section{Existência e Estabilidade}

Lorem ipsum dolor sit amet, consectetur adipiscing elit. Sed do eiusmod tempor incididunt ut labore et dolore magna aliqua.

Ut enim ad minim veniam, quis nostrud exercitation ullamco laboris nisi ut aliquip ex ea commodo consequat.

\section{Reflexões e Discussão}

Lorem ipsum dolor sit amet, consectetur adipiscing elit. Sed do eiusmod tempor incididunt ut labore et dolore magna aliqua.

Ut enim ad minim veniam, quis nostrud exercitation ullamco laboris nisi ut aliquip ex ea commodo consequat.

\section{Considerações Finais}

Lorem ipsum dolor sit amet, consectetur adipiscing elit. Sed do eiusmod tempor incididunt ut labore et dolore magna aliqua.

Ut enim ad minim veniam, quis nostrud exercitation ullamco laboris nisi ut aliquip ex ea commodo consequat.


\nocite{*}
\bibliographystyle{abntex2-alf}
\bibliography{library}

\end{document}
